
\subsubsection*{1.6.1}
\phantomsection
\label{sga3:VIIB:1.6.1}

Let \(Y\) be an étale formal \(k\)-variety and let
\(f:X\to Y\) be a morphism in \(\operatorname{Vaf}/k\). Then the
following square is Cartesian, where \(\Gamma_f\) is the graph
morphism \(X\to X\times Y\), with components
\(\operatorname{id}_X\) and \(f\):

\begin{center}
  \phantomsection
  \label{sga3:VIIB:1.6.1:cartesian-square}
  \begin{tikzcd}[row sep=2.8em,column sep=3.7em]
X \arrow[r,"\Gamma_f"] \arrow[d,"f"']
  & X\times Y \arrow[d,"{f\times\operatorname{id}_Y}"] \\
Y \arrow[r,"\Delta_Y"]
  & Y\times Y\, .
\end{tikzcd}

\end{center}

It follows that \(\Gamma_f\) is a local isomorphism, and hence that
\[
  f=\operatorname{pr}_Y\circ\Gamma_f
\]
is topologically flat if \(\operatorname{pr}_Y\) is; this holds, for
example, if \(X\) is topologically flat over \(k\).

Conversely, since \(Y\to k\) is topologically flat, \(X\to k\) will
also be topologically flat if \(f\) is \textup{(cf.\ \SGARefLink{sga3:VIIB:1.3.3}{1.3.3})}. Taking
in particular for \(f\) the canonical morphism
\[
  p:X\longrightarrow X_e
\]
of \SGARefLink{sga3:VIIB:1.6}{1.6}, we obtain the following.

\medskip
\noindent\textbf{Proposition.}
\phantomsection
\label{sga3:VIIB:1.6.1.prop}
{\itshape Let \(X\) be a formal variety over \(k\). The morphism
\(X\to X_e\) is topologically flat if and only if \(X\) is
topologically flat over \(k\).}

\medskip
\noindent\textbf{Remark 1.6.2.}\footnote{%
\label{sga3:VIIB:nde:74}The editors inserted this remark here; in the
original it appeared in \SGARefLink{sga3:VIIB:2.5.2:proposition}{2.5.2}.}
\phantomsection
\label{sga3:VIIB:1.6.2}
When \(k\) is a \emph{perfect} field, the functor
\[
  X\longmapsto X_e
\]
is also right adjoint to the inclusion
\[
  \operatorname{Vaf}^{\mathrm{\acute{e}t}}/k
  \lhook\joinrel\longrightarrow
  \operatorname{Vaf}/k.
\]
Indeed, in this case
\[
  \mathcal O_{X_e,x}=\kappa(x)
\]
for every \(x\in X\), and the canonical projections
\[
  \mathcal O_{X,x}\longrightarrow\kappa(x)
\]
define a morphism \(s:X_e\to X\), which is a section of
\(p:X\to X_e\). For every morphism \(f:X\to Y\), the following
diagrams are commutative:

\begin{center}
  \phantomsection
  \label{sga3:VIIB:1.6.2:local-ring-square}
  \begin{tikzcd}[row sep=2.6em,column sep=3.5em]
\mathcal O_{Y,f(x)} \arrow[r] \arrow[d]
  & \mathcal O_{X,x} \arrow[d] \\
\kappa(f(x)) \arrow[r]
  & \kappa(x)
\end{tikzcd}

  \hspace{.07\textwidth}
  \phantomsection
  \label{sga3:VIIB:1.6.2:functoriality-square}
  \begin{tikzcd}[row sep=2.6em,column sep=3.2em]
X \arrow[r,"f"] & Y \\
X_e \arrow[r,"f_e"] \arrow[u,"s_X"]
  & Y_e \arrow[u,"s_Y"']
\end{tikzcd}

\end{center}

Thus \(s\) is functorial in \(X\), and it follows that
\(X\mapsto X_e\) is right adjoint to the inclusion
\[
  \operatorname{Vaf}^{\mathrm{\acute{e}t}}/k
  \lhook\joinrel\longrightarrow
  \operatorname{Vaf}/k.
\]
Consequently, being a right adjoint, the functor \(X\mapsto X_e\)
commutes with projective limits whenever they exist in
\(\operatorname{Vaf}^{\mathrm{\acute{e}t}}/k\),\footnote{%
\label{sga3:VIIB:nde:75}This is the case for finite projective
limits.} and hence, in particular, with finite products.

\noindent
(This can also be verified directly: for every formal
\(k\)-variety \(X\), with affine algebra \(A\), the affine algebra of
\(X_e\) is the quotient of \(A\) by its radical \(r(A)\); and, since
\(k\) is perfect, the quotient of
\(\mathcal O_{X,x}\widehat{\otimes}\mathcal O_{Y,y}\) by its radical
is the algebra \(\kappa(x)\otimes_k\kappa(y)\), because the latter is
semisimple.)

